\documentclass[12pt]{article}
\usepackage[a3paper,landscape,margin=0pt]{geometry}
\usepackage{tikz}
\usepackage{xcolor}
\usepackage{graphicx}
\usepackage[fontset=fandol]{ctex}
\usetikzlibrary{positioning,calc,fadings,decorations.pathmorphing,shadows}

% 优先用 PNG（已转换好的），fallback 到 jpg 和绝对路径
\graphicspath{{images_png/}{images/}{/home/ubuntu/Projects/poster4/images_png/}}

\definecolor{deepred}{HTML}{B01B22}
\definecolor{gold}{HTML}{D4AF37}
\definecolor{darkbg}{HTML}{1A1A1A}

% 新增第二个背景所需的颜色
\definecolor{bgdeepred}{RGB}{139,0,0}
\definecolor{bggold}{RGB}{255,215,0}
\definecolor{bgbrightgold}{RGB}{255,240,150}

\pagestyle{empty}

\begin{document}
	\begin{tikzpicture}[remember picture, overlay]

		% ========== 背景（来自第二个文档的风格） ==========
		% 深红底色
		\fill[bgdeepred] (current page.south west) rectangle (current page.north east);

		% 同心圆
		\fill[red!50!black, opacity=0.55] (current page.center) circle (34);
		\fill[red!60!black, opacity=0.45] (current page.center) circle (26);
		\fill[red!70!black, opacity=0.35] (current page.center) circle (18);

		% 顶部横向条纹
		\foreach \i in {0,0.5,...,6} {
			\fill[red!75!black, opacity={0.05+0.06*\i}]
			($(current page.north west)+(0,-\i*1.2)$) rectangle ($(current page.north east)+(0,-\i*1.2-0.6)$);
		}

		% 顶部光晕
		\fill[red!95!black, opacity=0.45] (current page.north) circle (23);
		\fill[red!98!black, opacity=0.35] (current page.north) circle (15);

		% 放射光芒
		\foreach \angle in {0,15,...,345} {
			\draw[bggold!25!red, line width=0.6pt, opacity=0.35]
			(current page.center) -- ++(\angle:47);
		}
		\foreach \angle in {7.5,22.5,...,352.5} {
			\draw[bggold!40!red, line width=1.2pt, opacity=0.18]
			(current page.center) -- ++(\angle:42);
		}
		\foreach \angle in {200,210,...,340} {
			\draw[bgbrightgold, line width=0.5pt, opacity=0.25]
			(current page.north) -- ++(\angle:31);
		}

		% 金色粒子
		\foreach \x/\y/\r in {
			5.78/27.84/0.09, 11.81/26.98/0.08, 30.19/27.47/0.11, 36.23/26.61/0.08,
			3.94/20.42/0.08, 38.06/20.79/0.09, 7.35/15.47/0.06, 34.65/15.10/0.08,
			4.73/9.28/0.08, 37.28/9.65/0.06, 8.40/11.76/0.06, 33.60/11.39/0.08,
			14.44/22.89/0.08, 27.56/22.52/0.06, 17.85/28.22/0.08, 24.15/27.84/0.09,
			11.03/3.96/0.06, 30.98/4.33/0.08, 6.56/6.44/0.06, 35.44/6.81/0.08,
			9.98/17.57/0.06, 32.03/17.94/0.08, 13.65/12.99/0.06, 28.35/13.37/0.06
		} {
			\fill[bgbrightgold, opacity=0.85] (\x,\y) circle (\r);
		}

		% ========== 顶部标题 ==========
		\node[anchor=north] at ($(current page.north)+(0,-1.0)$) {
			\begin{tikzpicture}
				\node[fill=deepred, rounded corners=4pt, inner sep=14pt,
				drop shadow={shadow xshift=3pt, shadow yshift=-3pt, black, opacity=0.6}]
				(titlebox) {
					\sffamily\bfseries\color{white}\fontsize{60}{68}\selectfont
					建军百年 · 强军有我
				};
				\node[below=6pt of titlebox, text=gold,
				font=\sffamily\fontsize{22}{26}\selectfont] {
					增强国防观念 \quad|\quad 履行国防义务 \quad|\quad 共筑钢铁长城
				};
			\end{tikzpicture}
		};

		% ========== 中部六图矩阵（PNG） ==========
		\coordinate (gridcenter) at ($(current page.center)+(0,-1.2)$);

		\foreach \file/\cap/\xpos/\ypos in {%
			{tank.png}/{陆战铁甲}/{-9.1}/{3.4},%
			{missile.png}/{大国长剑}/{0}/{3.4},%
			{ship.png}/{向海图强}/{9.1}/{3.4},%
			{jet.png}/{制空雄鹰}/{-9.1}/{-3.4},%
			{emblem.png}/{军魂永驻}/{0}/{-3.4},%
			{logo.png}/{忠诚担当}/{9.1}/{-3.4}%
		}{%
			\begin{scope}[shift={($(gridcenter)+(\xpos cm,\ypos cm)$)}]
				\node[inner sep=0pt] (pic) {%
					\includegraphics[width=8.5cm,height=6cm,keepaspectratio]{\file}%
				};
				\draw[gold, line width=1.5pt, rounded corners=2pt]
				(pic.south west) rectangle (pic.north east);
				\node[anchor=south, fill=deepred, text=white, opacity=0.92,
				text opacity=1, minimum width=8.5cm, inner sep=5pt,
				font=\sffamily\bfseries\fontsize{20}{24}\selectfont]
				at (pic.south) {\cap};
			\end{scope}
		}

		% ========== 底部标语 ==========
		\node[anchor=south] at ($(current page.south)+(0,1.0)$) {
			\begin{tikzpicture}
				\node[fill=gold, text=darkbg, rounded corners=3pt, inner sep=10pt,
				font=\sffamily\bfseries\fontsize{26}{30}\selectfont] {
					强国必须强军 \quad 军强才能国安
				};
			\end{tikzpicture}
		};
		% ========== 右下角签名 ==========
		\node[anchor=south east, text=gold,
		font=\sffamily\fontsize{10}{10}\selectfont]
		at ($(current page.south east)+(-1.8cm,1.2cm)$) {
			8A619
		};
		% ========== 四角装饰 ==========
		\foreach \corner/\rx/\ry in {
			south west/0/0, south east/1/0, north west/0/1, north east/1/1
		}{
			\begin{scope}[shift={(current page.\corner)}]
				\pgfmathsetmacro\sx{\rx==0?1:-1}
				\pgfmathsetmacro\sy{\ry==0?1:-1}
				\draw[gold, line width=3pt]
				(0.4*\sx,0.4*\sy) -- (2.2*\sx,0.4*\sy) -- (2.2*\sx,1.4*\sy);
				\draw[deepred, line width=1.5pt]
				(0.7*\sx,0.7*\sy) -- (1.6*\sx,0.7*\sy) -- (1.6*\sx,1.1*\sy);
			\end{scope}
		}

	\end{tikzpicture}
\end{document}